\title{Large Language Models Can Automatically Engineer Features for Few-Shot Tabular Learning}

\begin{abstract}
Large Language Models (LLMs) have demonstrated exceptional proficiency in addressing complex, novel reasoning tasks, presenting a valuable opportunity for improving tabular learning processes, which are integral across numerous real-world settings. This work introduces a new in-context learning paradigm, \textsf{FeatLLM}, that leverages LLMs for feature engineering to create an enhanced representation of the original dataset, optimized for tasks involving tabular predictions. The features generated by this approach are subsequently processed by a straightforward downstream machine learning model, such as linear regression, yielding robust performance even in few-shot learning contexts. Importantly, \textsf{FeatLLM} relies solely on this basic predictive model and the newly engineered features at inference, avoiding repeated dependence on the LLM. In contrast to existing LLM-driven methods, our framework dispenses with the requirement of querying the LLM for individual samples during inference, needs only API-level LLM access, and is not constrained by prompt length limitations. Empirical evaluation across a variety of tabular datasets and domains confirms that \textsf{FeatLLM} consistently produces superior, high-quality rules, achieving an average performance boost of 10\% over methods like TabLLM and STUNT.
\end{abstract}

\section{Introduction}
Large Language Models (LLMs) have, in recent times, exhibited impressive performance in generating relevant outputs for unfamiliar tasks, often circumventing the necessity for specific fine-tuning strategies~\cite{creswell2022selection,imani2023mathprompter,taylor2022galactica}. Through their extensive pretraining on massive text datasets, LLMs amass a breadth of prior knowledge that can, in many cases, obviate the reliance on specialized domain experts~\cite{Genomes2021,singhal2023large,wilcox2003role}. This adaptability positions them as pivotal tools for automating a wide array of applications, such as dialogue analysis~\cite{finch2023leveraging} and program repair~\cite{fan2023automated}. By supplying prompts that include descriptive task information and a small number of relevant examples, LLMs demonstrate an ability to capture task context and focus on pertinent features or sample cases~\cite{brown2020language,zhao2023survey}, reflecting their strengths in data analysis and suggesting their utility as advanced feature engineering agents.

The extension of LLMs’ reasoning abilities and embedded knowledge to tabular learning has been the focus of increasing attention. For example, leveraging prior insights—such as associations between age and disease risk—can enrich predictive modeling for health outcomes. Existing literature often represents both feature and task descriptions as natural language strings, restructures tabular data accordingly, and presents this serialized input to LLMs~\cite{dinh2022lift,wang2023anypredict}. Such efforts commonly utilize in-context learning~\cite{nam2023semi} or minimal-parameter fine-tuning techniques~\cite{dinh2022lift,hegselmann2023tabllm}. This incorporation of LLM-driven background knowledge has been shown to outperform classic tabular learning methods, especially when training data is scarce~\cite{hegselmann2023tabllm}.

Prevailing LLM-based techniques for tabular data analysis face several notable drawbacks. For prediction tasks conducted in an end-to-end manner, these methods necessitate performing a separate inference on the LLM for each data instance, thereby imposing substantial computational overhead. Achieving strong predictive accuracy typically depends on fine-tuning the underlying LLM, a strategy that proves impractical when using LLMs lacking publicly available training code or APIs for model adaptation. This is particularly problematic given that recent high-performing LLMs restrict access exclusively to inference APIs~\cite{anil2023palm,openai2023gpt4}. Another limitation is the poor scalability with extended prompts~\cite{liu2023lost}; as the number of attributes in the tabular data increases and the serialized representation lengthens, these methods frequently either break down or suffer significant drops in predictive quality. Collectively, these challenges create significant obstacles for the widespread adoption of LLM-based tabular learning solutions in applied settings.

Our method takes an alternative perspective by redefining the function of LLMs from end-to-end predictors to agents of feature engineering, capitalizing on their potential for effective knowledge reuse. We present a new in-context learning paradigm, \textsf{FeatLLM}, designed to reveal the implicit “criteria” underpinning LLM decision-making processes. As an example, for disease prediction, the LLM is tasked with inferring and formulating explicit rules that characterize the relationship between input features and disease outcomes. Such generated criteria, in the form of rules, are then employed to construct novel features which can supplant original variables during subsequent modeling stages. By reallocating the LLM's capabilities to feature synthesis, this process simplifies the requirements for downstream modeling; after feature extraction, there is no longer a need for complex inference-time models, thereby enhancing the speed of prediction. Moreover, this approach exploits the few-shot learning strengths of LLMs: by integrating appropriate demonstrations in the input prompts, relevant features can be harvested without any explicit model retraining. To further address limitations posed by long prompts, ensemble-inspired strategies such as feature bagging~\cite{breiman2001random} can be adopted, providing practical ways to manage and reduce prompt length.

\textsf{FeatLLM} organizes the process of designing novel features from prompts by decomposing it into two key sub-tasks: (1) comprehending the task and deducing how features relate to the targets; and (2) leveraging insights from the previous inference along with several exemplar cases to generate decisive discriminative rules. This methodical reasoning prompts LLMs to disregard features that do not contribute relevant information when constructing rules. The inferred rules are subsequently evaluated on the dataset, with the rules being programmatically encoded; the resulting data are recast as binary variables indicating if each sample satisfies the respective rules. A simple model is then employed to predict class probabilities using these binary features. The approach further enhances reliability through an ensemble strategy, wherein new features are repeatedly synthesized using feature bagging. By tailoring the counts of features and samples selected during this process, the issue of unwieldy prompt length is addressed. \textsf{FeatLLM} imposes no requirements for additional training and offers efficient inference, making LLM deployment practical in diverse applications.

\textsf{FeatLLM}’s efficacy is assessed on 13 tabular datasets within low-data scenarios, demonstrating its strong and stable results. Our experiments highlight that LLMs can effectively combine prior knowledge with data-driven information to generate features of high utility. Across multiple configurations, our method surpasses leading few-shot learning approaches. An anonymized GitHub repository hosts our code at~\texttt{https://github.com/Sungwon-Han/FeatLLM}.

\section{Related Work}

\subsection{Few-Shot Learning with Tabular Data} Few-shot learning has seen rapid progress, enabling the development of generalizable data representations despite very limited labeled examples and thus reducing annotation efforts~\cite{chen2018closer}. While its early successes centered on image-based tasks, few-shot paradigms have now been shown to transfer well to modalities such as text, audio, and more recently, tabular datasets~\cite{majumder2022few,nam2023stunt,schick2021exploiting}. One of the inherent challenges when learning from scant labeled data is the possibility of models capturing spurious patterns that do not generalize~\cite{bartlett2021deep}. To counteract this, researchers have begun to supplement training with external information sources or explicitly embed domain expertise to inject suitable inductive biases. This might include leveraging large, accessible unlabeled tabular collections together with purposefully crafted data augmentations under semi-supervised frameworks~\cite{ucar2021subtab,yoon2020vime}, or initializing models via broad, unsupervised pre-training not specific to any one task~\cite{somepalli2021saint}. Common unsupervised techniques involve tasks such as masked reconstruction—where parts of the input are hidden or perturbed and then reconstructed~\cite{arik2021tabnet,majmundar2022met}—or employing augmented data to enact contrastive learning schemes~\cite{bahri2022scarf,chen2020simple}. However, the inherent variability in tabular datasets complicates the design of universally effective augmentations, and these methods generally have yet to exhibit robust improvement in few-shot regimes~\cite{nam2023stunt}.

To address these issues, a line of recent research explores cross-dataset pre-training—training models on diverse tabular data sources not strictly tied to the downstream task, followed by transfer to the specific target context~\cite{zhang2023generative,levin2022transfer,zhu2023xtab}. For instance, TransTab~\cite{wang2022transtab} advances transformer models that encode semantic associations between table columns by utilizing column names alongside values. This accumulation of knowledge from heterogeneous tables has allowed for competent zero-shot and few-shot adaptation on new tasks. Conversely, TabPFN~\cite{hollmann2022tabpfn} introduces a Bayesian neural network pre-trained with synthetically generated datasets sampled from assorted distributions, facilitating direct inference on new tasks with access to only support and query samples and without additional fine-tuning. Through the integration of experiences from multiple sources, these approaches have consistently outperformed classic tree-based methods, particularly in data-scarce scenarios.

\subsection{Language-Interfaced Tabular Learning} Integrating the substantial prior knowledge embedded in large language models (LLMs)~\cite{anil2023palm,openai2023gpt4} presents another avenue for advancing tabular learning. LLMs, benefiting from training on immense and varied text datasets, have displayed strong capability to generalize to novel tasks, a property that has been leveraged in various application domains~\cite{brown2020language}. More recently, strategies have emerged that exploit the encyclopedic priors within LLMs by converting tabular data into textual representations and supplying these, along with feature and task descriptions, directly to the model as input prompts~\cite{dinh2022lift,hegselmann2023tabllm,wang2023anypredict}. This structure allows the model to interpret the relationships inherent in the data and task, facilitating informed predictions. Further enhancements involve either tailoring LLMs through lightweight parameter adaptation mechanisms—such as LoRA~\cite{hu2021lora} or IA3~\cite{liu2022few}—or leveraging in-context learning, where few-shot demonstrations are embedded within the prompt itself, eschewing any model weight updates altogether~\cite{dinh2022lift,hegselmann2023tabllm,nam2023semi,slack2023tablet}.

While the previously discussed approaches have significantly advanced few-shot tabular learning, their reliance on routing inference through large language models (LLMs) introduces practical obstacles in industrial settings. This work seeks to move away from treating the LLM as an end-to-end predictive black box for each individual sample. Rather than producing predictions for each sample directly, we focus the LLM on deriving the underlying rules or decision criteria corresponding to each class. \textsf{FeatLLM} subsequently translates these synthesized rules into programmatic feature generators. This enables downstream prediction tasks to operate independently of the LLM by utilizing these extracted rules, effectively harnessing in-context learning for rule induction based on few-shot examples and previously acquired knowledge.

\section{Methods}
The core functionality of \textsf{FeatLLM} is to obtain “rules”—sets of class-conditional conditions—by analyzing a small number of instances per class, as illustrated in Figure~\ref{fig:main_model}. These rules, produced via the LLM, are translated into binary features that represent, for any given sample, whether the sample fulfills a particular rule. The resulting suite of binary indicators is fed into a weighted linear (dot product) layer where all weights are constrained to be non-negative, yielding scores interpretable as class probabilities. Model training thus emphasizes learning which generated rule-features most strongly influence class likelihoods, as elaborated in Section~\ref{Sec:training}. To bolster robustness, this methodology is iteratively performed with bagging and the outputs are consolidated through an ensemble approach. A more granular overview of the problem definition and methodological steps follows.

\vspace*{-0.12in}
\paragraph{Problem formulation.} Consider a tabular dataset $\mathcal{D} = \{(\mathbf{x}^i, \mathbf{y}^i)\}_{i=1}^{N}$, which contains $N$ labeled data points. Each input $\mathbf{x}^i$ is represented as a $d$-dimensional vector corresponding to $d$ different features, and the dataset includes a set of feature names expressed in natural language, $F = \{f_j\}_{j=1}^{d}$, with their separate definitions. Under a classification framework, each label $\mathbf{y}^i$ takes a value from a predetermined collection of classes. For $k$-shot learning scenarios, the model is trained using only $k$ samples (with $k < N$), which are drawn randomly from the labeled dataset.

\subsection{Prompt Design for Extracting Rules}
\label{Sec:prompt_design}
To enhance the LLM's ability to derive rules by following sound reasoning steps, we structure the prompt to resemble the analytic approach of a human expert when addressing tabular prediction tasks. The constructed prompt consists of three primary sections (illustrated in Fig.~\ref{fig:prompt_for_rule} and detailed in Appendix~\ref{sec:example_rule}):

\vspace*{-0.12in}
\paragraph{Basic information description.} This segment supplies all foundational information necessary for task resolution (highlighted in orange in Fig.~\ref{fig:prompt_for_rule}). It comprises a statement outlining the problem to be addressed and its associated label information, summaries of the input features, and representative instances. The task description poses a specific query (for example, ``Does this patient's myocardial infarction complications data indicate chronic heart failure? Yes or no?"). Additional task examples are compiled in Appendix~\ref{sec:task_desc}. Feature descriptions specify the type of value each feature can take and may include further elaboration. For features that are categorical, sample category values can be listed. To illustrate, a selected few training examples, accompanied by their correct labels, are linearized into text according to previously established approaches~\cite{dinh2022lift,hegselmann2023tabllm}:

\begin{align}
&\text{Serialize}(\mathbf{x}^i,\mathbf{y}^i, F) = \nonumber \\ 
&``\ f_1\ \text{is}\ \mathbf{x}^i_1.\ ... \ f_d\  \text{is}\  \mathbf{x}^i_d.\ \ \text{Answer:}\  \mathbf{y}^i\ ”
\end{align}

\vspace*{-0.12in}
\paragraph{Reasoning instruction.} Our objective is to strengthen the reasoning capabilities of the LLM by supplying it with targeted reasoning instructions (highlighted in blue in Fig.~\ref{fig:prompt_for_rule}). This guidance commences with an opening statement akin to the chain-of-thought methodology~\cite{wei2022chain}. We then structure the inference task for the LLM into two phases. In the initial phase, the LLM is tasked with determining the causal associations or likely relationships between various features and the overall task description. Importantly, this step is performed independently of the provided example demonstrations, thereby prompting the LLM to depend on its inherent knowledge base or commonsense reasoning and to comprehensively leverage its existing understanding of the domain. The subsequent phase instructs the LLM to synthesize the information extracted in the first phase with the example demonstrations, and then to derive rules corresponding to each class. This bifurcated reasoning procedure helps in mitigating the influence of spurious correlations stemming from non-informative columns and directs the model's focus towards features of greater relevance. To avoid either an inadequate or excessive extraction of rules, which could result in underfitting or an impractically large prompt, we prescribe a fixed rule count (specifically, 10 rules)\footnote{For detailed analysis regarding the choice of rule numbers, refer to Section~\ref{sec:analysis}.} per class. Moreover, during the rule generation stage, we prompt the LLM to consult the feature type information presented within the basic description section to reduce the risk of introducing type incompatibility into the rules.

\vspace*{-0.12in}
\paragraph{Response instruction.} To support more effective interpretation and downstream use of the extracted rules, we direct the LLM on how to articulate its responses by providing concrete response guidelines (highlighted in yellow in Fig.~\ref{fig:prompt_for_rule}). The prompt delineates specific formatting requirements for the responses associated with each answer class (i.e., structure for the response) and specifies the arrangement each individual rule should follow (i.e., structure for the rule). Through these directions, the LLM is able to align the response format to the corresponding feature types as needed. Absent explicit instructions to the contrary, the LLM may conjoin multiple rules using logical connectors such as AND or OR. An illustrative example is furnished in Appendix~\ref{sec:outcome-rule}.

\subsection{Inferring Class Likelihood via Rules}
\label{Sec:training}

\paragraph{Parsing rules for feature generation.} The rules identified in the earlier step are used to construct novel binary features. These features, generated for every class, serve to indicate whether a particular sample fulfills the rule conditions relevant to that class (yielding a value of '0' or '1'). These binary indicators inform the prediction of the probability that a sample falls within a particular class. Since the rules are expressed in natural language by the LLM, developing a strategy for parsing these rules to enable automatic feature creation becomes essential. While we impose constraints on the structure and formatting of rules to aid in their later extraction, outputs from the LLM may nonetheless exhibit variable syntax and minor inconsistencies~\cite{kovavc2023large}. As an example, the LLM could reformulate mathematical expressions, writing ``greater than'' instead of ``$>$'' or ``less than'' instead of ``$<$,'' introducing extra variability that complicates direct parsing.

To circumvent the complexities of engineering bespoke parsing programs, we take advantage of the LLM’s strong capabilities in semantic interpretation and code synthesis (having been exposed to code during training). Instead of hand-crafting intricate logic, we incorporate the function signature, alongside input/output specifications and the target rules, in the prompt submitted to the LLM. As shown in Figures~\ref{fig:prompt_for_parsing} and ~\ref{sec:example_parsing}-\ref{sec:example_parse_outcome} in the Appendix, this approach is illustrated through several example prompts and their corresponding outputs. The resulting code is then run with Python’s \texttt{exec()} function to facilitate transformation of the data.

\vspace*{-0.12in}
\paragraph{Inferring class likelihood.} After producing new binary rule-based features for each class, a straightforward approach for estimating how likely a sample belongs to a particular class is to simply tally the number of satisfied rules for that class (that is, summing the binary indicators associated with each class). However, it is important to recognize that the contributions of different rules may vary, motivating the need to assess their relative importance using the training data. To accomplish this, we train a standard machine learning (ML) model on the binary features for each class. For instance, we can employ a bias-free linear model. Define $\mathbf{z}_k^i \in \{0, 1\}^R$ as the vector of binary features produced for class $k$ from a given sample $\mathbf{x}^i$, where $R$ is the total number of rules per class, and $c$ denotes the total number of classes. The predicted probability for each class $\mathbf{p}^i$ can then be determined as follows:

\begin{align}
\text{logit}_k^i &= \max(\mathbf{w}_k, 0) \cdot \mathbf{z}_k^i, \nonumber \\
\mathbf{p}^i &= \text{Softmax}({[\text{logit}_{1}^i, ..., \text{logit}_{c}^i]}). \label{eq:infer_prob}
\end{align}

In this context, $\mathbf{w}_k$ denotes the adjustable parameter in the linear framework, reflecting the relevance of each individual feature. By constraining these weights to reside within a positive domain, it is guaranteed that no feature produced by the LLM will detract from a class's likelihood during the training process. The resulting logit vector is subsequently transformed into class probabilities by applying the softmax function. The objective is to learn the optimal $\mathbf{w}_k$ via cross-entropy minimization, utilizing only a limited number of annotated examples for training.

\vspace*{-0.12in}
\paragraph{Ensembling with bagging.} To conclude, the entire inference procedure is executed multiple times to build a set of inference models, whose outputs are ultimately aggregated in an ensemble to yield final predictions. Specifically, after $T$ independent runs, each yielding learned parameters $\{\mathbf{w}[t]\}_{t=1}^T$, predictions are generated by taking the mean of the class probability vectors $\mathbf{p}^i$ computed from each $\mathbf{w}[t]$, as described in Eq.~\ref{eq:infer_prob}.

Randomness is injected into the rule generation process for ensembling by specifying a sufficiently elevated temperature during LLM querying. In addition, the few-shot examples are shuffled prior to inference because, as prior work has established, the sequencing of demonstrations can influence LLM responses~\cite{lu2022fantastically}\footnote{Further discussion on the diversity of rules can be found in Appendix~\ref{sec:diversity}}. To harness this variance for stronger predictive performance, bagging is used to subsample features or data instances for each trial, a strategy that proves increasingly advantageous as the training set or feature space grows. This ensemble scheme introduces two significant benefits. Firstly, occasional faulty rules from some LLM trials can be offset by correct rules from others, making the entire pipeline more robust. This leverages the LLM's property of self-consistency~\cite{wang2022self}: rules that frequently appear across runs are more likely to be reliable. Secondly, ensembling mitigates the challenge posed by LLM prompt size limitations. If the tabular input surpasses permissible prompt length, bagging curtails the input size, safeguarding the model’s reliability. While a single rule set may fail to capture interactions among all variables, aggregating predictions from several diverse, weak learners across multiple trials counteracts the potential drawbacks of partial coverage in individual trials.

\section{Experiments}
We assess the performance of \textsf{FeatLLM} across a range of tabular datasets and carry out comprehensive analyses to better understand the workings of our framework. Owing to space limitations, supplemental experiments—including variations in LLM selection, qualitative studies, and others—are presented in the Appendix.

\subsection{Performance Evaluation}
\paragraph{Datasets.}
A total of 13 tabular datasets are utilized in our empirical study, focusing on binary and multi-class classification settings: (1) Adult~\cite{asuncion2007uci} for predicting if an individual's annual income surpasses \$50,000; (2) Bank~\cite{moro2014data} for classifying whether a client will subscribe to a term deposit; (3) Blood~\cite{yeh2009knowledge}, targeting the likelihood of donors returning for future donations; (4) Car~\cite{kadra2021well} for estimating car quality categories; (5) Communities~\cite{misc_communities_and_crime_183} for predicting the regional crime rate; (6) Credit-g~\cite{kadra2021well}, used to assess whether an individual's credit risk is favorable or unfavorable; (7) Diabetes\footnote{\texttt{kaggle.com/datasets/uciml/pima-indians-diabetes-database}} to determine the presence of diabetes; (8) Heart\footnote{\texttt{kaggle.com/datasets/fedesoriano/heart-failure-prediction}} to forecast occurrences of coronary artery disease; and (9) Myocardial~\cite{misc_myocardial_infarction_complications_579}, aimed at detecting chronic heart failure.

To broaden the evaluation, we include newly released datasets as well as synthetic benchmarks that have not been widely studied and are unlikely to have been seen during LLM pre-training: (10) Cultivars, which predicts soybean cultivar grain yield levels\footnote{\texttt{archive.ics.uci.edu/dataset/913}}; (11) NHANES, tasked with inferring a person’s age group based on the record\footnote{\texttt{archive.ics.uci.edu/dataset/887}}; (12) Sequence-type, responsible for labeling sequences as either arithmetic, geometric, Fibonacci, or Collatz; and (13) Solution-mix, which identifies whether a four-component mixture exceeds a concentration of 0.5. The first two among these were only incorporated into the UCI Machine Learning Repository after September 2023, thereby guaranteeing exclusion from pre-training data of the LLM API (gpt-3.5-turbo-0613) applied by our method. The last two are custom-designed synthetic datasets, precluding the possibility of memorization by the LLM API and thus supporting a fair assessment.

The datasets span a spectrum of scales and intricacy, as documented in Table~\ref{Tab:basic-desc}. Each one features explicitly defined attribute names and explanations. In particular, datasets such as `Communities' and `Myocardial' involve upwards of 100 attributes, thus straining the limited context capacity inherent to most LLMs.

\vspace*{-0.12in}
\paragraph{Baselines.}
We benchmark against nine baseline methods. The initial three are standard tabular data models: (1) Logistic Regression (LogReg), (2) XGBoost~\cite{chen2016xgboost}, and (3) RandomForest~\cite{ho1995random}. The subsequent two operate under the assumption of access to ample unlabeled data, including (4) SCARF~\cite{bahri2022scarf} and (5) STUNT~\cite{nam2023stunt}; however, procuring extensive unlabeled datasets is often not viable in practice. Another state-of-the-art approach is (6) TabPFN~\cite{hollmann2022tabpfn}, which leverages pretraining via synthetically created data from diverse distributions. Lastly, three LLM-based techniques are evaluated: (7) In-context learning (In-context)~\cite{wei2022emergent}, (8) TABLET~\cite{slack2023tablet}, and (9) TabLLM~\cite{hegselmann2023tabllm}. The in-context learning paradigm inserts a limited set of labeled instances directly into the input prompt, foregoing model fine-tuning. TABLET enhances prompts through the integration of supplementary hints produced by an external classifier via rule sets and prototypes, aiming to improve predictive accuracy. TabLLM adopts efficient fine-tuning strategies such as IA3~\cite{liu2022few}, enabling few-shot tuning of the LLM for classification tasks.

\vspace*{-0.12in}
\paragraph{Implementation details.}
\textsf{FeatLLM} adopts GPT-3.5 as its primary LLM backbone, while maintaining compatibility with alternative LLM architectures (refer to Appendix~\ref{sec:backbones} for comparative backbone experiments). During LLM inference, the temperature parameter is configured at 0.5 and top-p remains at the API's default value of 1. For ensemble generation and rule extraction, we employ 20 ensembles and 10 rules, respectively. Comprehensive discussions of hyper-parameter effects are available in Figure~\ref{fig:hyperparameter2} and Appendix~\ref{sec:hyper-parameter}.

A linear model is optimized with the Adam algorithm, utilizing a learning rate of 0.01 across 200 epochs. Model selection for the best epoch relies on $k$-fold cross-validation. In our baseline implementations, GPT-3.5 serves as the LLM for in-context learning, while T0~\cite{sanh2022multitask} is utilized with TabLLM to align with its original setup. It is important to note that TabLLM only supports LLMs with publicly available checkpoints, precluding the use of GPT-3.5 for this baseline. For classical machine learning approaches, including LogReg, XGBoost, and RandomForest, hyper-parameters were optimized via grid search in tandem with $k$-fold cross-validation. The value of $k$ is set to either 2 or 4, ensuring that all classes are represented in the training set by at least one instance. Further implementation specifics can be found in Appendix~\ref{sec:baselines}.

\vspace*{-0.12in}
\paragraph{Results.}
Tables~\ref{tab:main1} and \ref{tab:main2} provide a detailed comparison of performance metrics across multiple datasets. To ensure the robustness of our findings, each experiment was conducted three times using different random splits of the training and test data; we report the mean and standard deviation of the resulting AUC scores. \textsf{FeatLLM} demonstrates superior performance, either outperforming or closely matching the top baselines in almost all scenarios. Remarkably, with only 4 shots, our framework achieves results that rival established tabular learning approaches that utilize 32 shots, as illustrated in Fig.~\ref{fig:compares}. While both STUNT and TabPFN showed notable gains for specific datasets, they did not offer consistent overall improvements over traditional machine learning baselines. Furthermore, LLM-based models such as in-context learning, TABLET, and TabLLM were unable to process datasets with a large number of features. Our methodology, integrating bagging and ensemble strategies, remains effective on high-dimensional datasets and consistently delivers robust performance. It should be highlighted that although our bagging and ensemble principles could be incorporated into current LLM-based frameworks, such integration would require doubling the inference operations for each sample, resulting in considerable computational overhead and making the approach less feasible in practice.

\subsection{Analysis \& Discussion}
\label{sec:analysis}
\paragraph{Ablation study.} We assess the relative impact of each major component on overall performance by conducting a set of ablation experiments, focusing specifically on: (1) \textsf{-Tuning}, in which the linear model’s weight tuning is skipped and instead logits are calculated by summing binary features; (2) \textsf{-Ensemble}, where the ensembling mechanism is omitted; (3) \textsf{-Description}, where feature descriptions are removed; and (4) \textsf{-Reasoning}, in which Step-1 of the rule generation instruction is excluded.

Table~\ref{tab:ablation} reports how these ablations alter the average AUC across all datasets. In every scenario, removal or modification of a component leads to diminished performance. Notably, the impact of weight tuning becomes increasingly pronounced as the number of available shots grows, suggesting that richer datasets enable better estimation of rule importance, amplifying the benefit of weight optimization. Conversely, eliminating feature descriptions and reasoning steps has a larger negative impact in few-shot settings, emphasizing the role of prior knowledge encoded in LLMs when data is limited. Finally, our ensemble-based strategy yields consistent performance improvements across all experimental conditions.

\vspace*{-0.12in}
\paragraph{Training \& inference time comparison.} We analyze the training and inference durations across several approaches. All experiments are conducted using the Adult dataset with a 16-shot training split. Unless otherwise noted, all models employ a single A100 GPU, with the exception of TabLLM, which utilizes four A100 GPUs to enable model parallelism. Table~\ref{tab:cost} summarizes the timing results. Importantly, \textsf{FeatLLM} achieves a notably low inference time, on par with standard machine learning models. In contrast, LLM-driven solutions like In-context learning, TABLET, and TabLLM are characterized by considerably longer inference times. With respect to training durations, \textsf{FeatLLM}'s efficiency aligns with that of other LLM-based strategies—it involves only 30 API invocations, a number that is budget-friendly and can be parallelized without significant overhead.

\vspace*{-0.12in}
\paragraph{Quality of features generated by \textsf{FeatLLM}.}
To assess the usefulness of the rule-based features generated via \textsf{FeatLLM} in practical settings, we perform an experiment focusing on their informativeness. Specifically, we compute the mean AUROC between these derived features and the ground-truth target values, followed by evaluating the proportion of features per dataset achieving an AUROC above 0.5. The outcomes are summarized in Table~\ref{tab:quality}. The data reveal that most generated rules are strongly related to the predictive target and deliver substantive information to the task at hand (see the column ``Before tuning"). Further, when fitting the linear model to the data, less informative rules—defined as those whose learned coefficients fall below a certain cutoff—are automatically excluded from the final model (see the column ``After tuning"). Here, the cutoff threshold is set to 0.1, as informed by the global distribution of the feature weights.

\vspace*{-0.12in}
\paragraph{Effect of adding spurious correlations.} To evaluate the capacity of different approaches to eliminate noisy spurious correlations under limited data conditions, we performed an investigation. Specifically, our experimental setup involved the Heart dataset, which we augmented by randomly appending columns derived from the Adult dataset; these additional columns have no relevance to the original Heart task. Model performances were subsequently assessed as the number of extraneous columns increased incrementally. To maintain consistency, all models were provided with the same 8 labeled instances and we did not consider any unlabeled data, thereby omitting approaches such as SCARF and STUNT from the scope of this analysis.

Performance variation arising from the introduction of spurious correlations is depicted in Figure~\ref{fig:spurious}. In contrast with standard baselines, \textsf{FeatLLM} experienced the smallest decline in performance when faced with these noisy features. This advantage likely stems from the mechanism within our framework that leverages prior domain knowledge to tie the relevance of features directly to the task during the process of rule extraction. As a result, the likelihood of the model deriving rules from unrelated attributes is diminished, thus enhancing robustness. Empirically, we found that only 1-2\% of generated rules involved noisy features. However, we also noted that if columns randomly appended are drawn from data semantically similar to the main task, \textsf{FeatLLM} remains susceptible to picking up spurious associations and misconstruing the causal connections between features and the task, as detailed further in Appendix~\ref{sec:spurious-appendix}.

\vspace*{-0.12in}
\paragraph{Hyper-parameter analysis}
\textsf{FeatLLM} involves two principal hyper-parameters: the ensemble count and the number of rules per class obtained from each LLM invocation. To assess their influence on the method's effectiveness, we performed systematic evaluations. Empirical evidence reveals that augmenting the ensemble size consistently elevates performance, but such improvements plateau past approximately 20 ensembles (refer to Fig.~\ref{fig:appendix-hyperparameter1} in Appendix). Regarding the number of rules, although variations in this parameter did not show statistical significance, we identified 10 as an optimal setting—balancing prompt length limitations against empirical performance (see Fig.~\ref{fig:hyperparameter2}). We hypothesize that including more rules expands the feature space, increasing the available information, but may induce overfitting, particularly in scenarios with few samples.

\section{Conclusion}
We introduce \textsf{FeatLLM}, which represents a significant advance in leveraging LLMs for few-shot tabular data analysis. Unlike previous approaches where LLMs make sample-wise predictions, \textsf{FeatLLM} utilizes these models for formulating predictive criteria akin to feature engineering. By combining rule generation with ensemble methods (bagging), \textsf{FeatLLM} operates with negligible inference overhead, obviates the need for model retraining, and works robustly even as feature counts grow. Empirical results demonstrate that \textsf{FeatLLM} delivers state-of-the-art predictive accuracy and offers a practical, resource-efficient solution for applying LLMs to real-world tabular problems.

Currently, \textsf{FeatLLM} is targeted toward few-shot settings. Looking ahead, we plan to broaden its capabilities to create new features in larger datasets, to handle diverse feature representations beyond traditional rule-based formats, and to enhance model interpretability—aiming to expand its applicability to a broader spectrum of tasks.

\section{Broader Impact}
The design of \textsf{FeatLLM} seeks to exploit the expansive knowledge embedded within LLMs, surpassing the performance ceilings of conventional supervised tabular methods, especially when labeled data is scarce. Our approach contributes to the broader goal of alleviating overfitting by leveraging such prior knowledge, advancing the efficiency of learning from limited data. \textsf{FeatLLM} stands to benefit practitioners facing labeling constraints—prevalent in domains like finance and healthcare, where annotation is costly or difficult—and can capitalize on the contextual expertise LLMs derive from large-scale, domain-relevant corpora. As deployment of this technique grows, it will be increasingly crucial to examine the effects of embedded biases and misinformation present in LLM pretraining sources, as this latent knowledge could influence or even overshadow the patterns observed in labeled datasets.

\end{document}